% !TeX program = xelatex
% !TeX encoding = UTF-8
% !BIB program = bibtex
% !TeX root = main.tex
\documentclass[conference]{IEEEtran}
\IEEEoverridecommandlockouts

% 中文支持（需 XeLaTeX 编译）：可用中文写草稿，定稿时替换为英文即可，无需改编译配置
\usepackage[UTF8, scheme=plain]{ctex}

\usepackage{cite}
\usepackage{amsmath,amssymb,amsfonts}
\usepackage{algorithm}
\usepackage{algorithmic}
\usepackage{graphicx}
\usepackage{textcomp}
\usepackage{xcolor}
\usepackage{booktabs}
\usepackage{multirow}
\usepackage{subfig}
\def\BibTeX{{\rm B\kern-.05em{\sc i\kern-.025em b}\kern-.08em T\kern-.1667em\lower.7ex\hbox{E}\kern-.125emX}}

\graphicspath{{figures/}}

\begin{document}

\title{Paper Title Here\\
% Sub-title should not be used for IEEE Xplore
\thanks{Identify applicable funding agency here. If none, delete this.}
}

\author{\IEEEauthorblockN{1st Given Name Surname}
\IEEEauthorblockA{\textit{Department of ...} \\
\textit{University of ...}\\
City, Country \\
email@example.com}
\and
\IEEEauthorblockN{2nd Given Name Surname}
\IEEEauthorblockA{\textit{Department of ...} \\
\textit{University of ...}\\
City, Country \\
email@example.com}
}

\maketitle

\begin{abstract}

Replace this placeholder with a concise summary of the problem, the proposed approach, and the key results. Do not use symbols, special characters, footnotes, or math in the abstract.
\end{abstract}

\begin{IEEEkeywords}
keyword one, keyword two, keyword three.
\end{IEEEkeywords}

\section{Introduction}
\IEEEPARstart{T}{his} opening paragraph provides the general background of the work: the application scenario, the driving trend, and why the problem is worth studying now. Keep it accessible to the general reader of the proceedings and quantify the impact with one or two illustrative figures of merit, then narrow down to the specific problem addressed in this paper.

The second paragraph reviews the state of the art at a high level and identifies the gap. Summarize the dominant existing approaches in one or two sentences each, explain their common limitation, and state explicitly what is still missing. This motivates the contributions of this paper, typically summarized as a bullet list:
\begin{itemize}
\item We formulate ... as ...
\item We propose a ... algorithm to ...
\item Theoretical analysis and simulation results demonstrate that ...
\end{itemize}

The closing paragraph of the introduction states the headline result in one or two sentences, e.g., ``the proposed scheme achieves a 20\% gain over the strongest baseline under typical operating conditions,'' which anchors the contribution list above.

\subsection{Related Work}
Review prior work and position this paper against existing contributions~\cite{ref_shannon1948, ref_boyd2004}. Organize the literature by methodology rather than chronologically: one short paragraph per family of approaches, contrasting their assumptions and performance, and closing with the sentence that separates this work from them.

A second paragraph may cover the complementary line of work in~\cite{ref_boyd2004} and explain why porting those tools to the present setting is non-trivial, e.g., new coupling constraints or a different information structure.

\subsection{Contributions and Organization}
Summarize the contributions and outline the remainder of the paper. Section~\ref{sec:model} describes the system model. Section~\ref{sec:method} presents the proposed method. Section~\ref{sec:results} reports the results, and Section~\ref{sec:conclusion} concludes the paper.

placeholder

placeholder

placeholder

placeholder

placeholder

placeholder

placeholder

placeholder

placeholder

placeholder

placeholder

placeholder

placeholder

placeholder

placeholder

placeholder

placeholder

placeholder

placeholder

placeholder

placeholder

placeholder

placeholder

placeholder

placeholder

placeholder

placeholder

placeholder

placeholder

placeholder

placeholder

placeholder

placeholder

placeholder

placeholder

placeholder

placeholder

\section{System Model}\label{sec:model}

\subsection{System Description}

We consider the multi-UAV mmWave ISAC network illustrated in Fig.~\ref{fig:system}. It consists of $K$ UAVs ($\mathcal K={1,2,\ldots,K}$ ), $M$  ground users ($\mathcal M$), $N$ static ground targets ($\mathcal N$), and a central controller (CC) which served by leader UAV.

The UAVs are equipped with $N_{t}$-element uniform linear arrays (ULAs) and fly at a fixed altitude band. Each UAV transmits a superposition of $M$ communication streams and one dedicated sensing stream, and it receives the target echoes of that transmission. Every ground user and target is considered static and equipped with a single antenna. The CC collects state information, runs the trajectory, power, and beamforming policy, and returns the resulting control to the UAVs. The mission horizon is divided into $T$ slots of duration $\Delta t$.

\subsection{A2G mmWave Channel Model}
\textbf{Probabilistic LoS model.} Let $\theta_{\rm el}$ denote the elevation angle (in degrees) of the UAV as seen from the ground node. The LoS probability follows the sigmoid model of Al-Hourani \emph{et al.}:
\begin{equation}
P_{\rm LoS}(\theta_{\rm el})=\frac{1}{1+a\exp\!\big(-b(\theta_{\rm el}-a)\big)}
\end{equation}
where $a,b$ are environment-dependent constants.

\textbf{Large-scale fading.} With $d_{k,m}=\max\{\|\mathbf p_k-\mathbf u_m\|,d_0\}$ and reference gain $\beta_0=(\lambda/4\pi d_0)^2$, the LoS/NLoS-weighted average path-loss gain is
\begin{equation}
\bar\beta_{k,m}=P_{\rm LoS}\,\beta_0\Big(\frac{d_0}{d_{k,m}}\Big)^{\alpha_{\rm L}} +\big(1-P_{\rm LoS}\big)\beta_0\Big(\frac{d_0}{d_{k,m}}\Big)^{\alpha_{\rm N}},
\end{equation}
where $\alpha_{\rm L}<\alpha_{\rm N}$ are the LoS and NLoS path-loss exponents.
 
\textbf{Small-scale fading.} The channel from UAV $k$ to user $m$ follows a Rician model:
\begin{equation}
\mathbf h_{k,m}=\sqrt{\frac{\kappa\bar\beta_{k,m}}{\kappa+1}}\, \mathbf a(\theta_{k,m},\phi_{k,m}) +\sqrt{\frac{\bar\beta_{k,m}}{\kappa+1}}\,\mathbf g_{k,m},
\end{equation}
with $\mathbf g_{k,m}\sim\mathcal{CN}(\mathbf 0,\mathbf I_{N_t})$ and Rician $K$-factor $\kappa$. 

ULA steering vector
\begin{equation}
[\mathbf a(\theta,\phi)]_n=\exp\!\Big(j2\pi\frac{d_a}{\lambda}\,n\cos\phi\cos\theta\Big), \quad n=0,\dots,N_t-1.
\end{equation}

\textbf{notes for steering vector.} 
The steering vector is \emph{not} power-normalized, i.e., $\|\mathbf a\|^2=N_t$, so that $\mathbb{E}\|\mathbf h_{k,m}\|^2=\bar\beta_{k,m}N_t$. The array gain is thus embedded in the channel, and $\|\mathbf w\|^2$ equals the actual transmit power of a beam.
 
\textbf{Remark.} We adopt an ``expected-channel'' approximation: the LoS/NLoS states are averaged in the large-scale gain before the Rician component is applied. A strict two-state model yields a bimodal received-power distribution, and by Jensen's inequality the present model slightly overestimates SNR-concave metrics such as the ergodic rate. We validate the approximation against the two-state model via Monte Carlo simulation.

\subsection{Transmit Signal and Downlink Communication}
UAV $k$ transmits
\begin{equation}
\mathbf x_k=\sum_{j=1}^{M}\mathbf w_{k,j}s_j+\mathbf w_{k,r}s_r,
\end{equation}
where $\mathbf w_{k,j}\in\mathbb C^{N_t}$ is the communication beamformer for user $j$, $\mathbf w_{k,r}$ is the dedicated sensing beamformer, $s_j$ is the data symbol of user $j$, and $s_r$ is a known probing waveform. The symbols satisfy $\mathbb{E}[s_js_i^*]=\delta_{ji}$, $\mathbb{E}[s_js_r^*]=0$, and $\mathbb{E}[|s_r|^2]=1$. Each UAV maintains one stream per user; streams of unserved users are driven to zero power, so the dimension of the decision space is independent of the association.

\textbf{Transmit covariance and power constraint.} Under the above symbol statistics,
\begin{equation}
\mathbf R_k=\mathbb{E}[\mathbf x_k\mathbf x_k^H]=\sum_{j=1}^{M}\mathbf w_{k,j}\mathbf w_{k,j}^H+\mathbf w_{k,r}\mathbf w_{k,r}^H\succeq\mathbf 0,
\end{equation}
\begin{equation}
\operatorname{tr}(\mathbf R_k)\le P_{\max}.
\end{equation}

\textbf{SINR and achievable rate.} The desired signal power of user $m$ is $S_m=|\mathbf h_{k(m),m}^H\mathbf w_{k(m),m}|^2$, where $k(m)$ denotes the serving UAV of user $m$. The interference comprises (i) multiuser interference (MUI) from the other streams of the serving UAV, (ii) inter-UAV interference from all communication streams of non-serving UAVs, and (iii) sensing-induced interference from the sensing beams of all UAVs:
\begin{multline}
I_m=\sum_{k\in\mathcal K}\sum_{j\in\mathcal M}\big|\mathbf h_{k,m}^H\mathbf w_{k,j}\big|^2\,\mathbf 1\{(k,j)\neq(k(m),m)\}\\
+\sum_{k\in\mathcal K}\big|\mathbf h_{k,m}^H\mathbf w_{k,r}\big|^2,
\end{multline}
\begin{equation}
\gamma_m=\frac{S_m}{\sigma^2+I_m},\qquad R_m=B\log_2(1+\gamma_m).
\end{equation}

\subsection{Sensing Model}
\textbf{Echo model.} The reflection coefficient of target $n$ with respect to UAV $k$ is
\begin{equation}
\alpha_{k,n}=\sqrt{\frac{\lambda^2\sigma_{\rm rcs}}{(4\pi)^3d_{k,n}^4}}\;e^{j\psi_{k,n}},
\end{equation}
which captures the radar cross-section (RCS), the two-way path loss $d^{-4}$, and an unknown phase $\psi_{k,n}$. After matched filtering with each known symbol and coherent integration over $N_s=BT_s$ samples, the observation vector stacked over the $M{+}1$ streams is
\begin{equation}
\mathbf y_{k,n}=\sqrt{N_s}\,\alpha_{k,n}\mathbf W_k^H\mathbf a(\theta_{k,n},\phi_{k,n})+\mathbf z,
\end{equation}
with $\mathbf z\sim\mathcal{CN}(\mathbf 0,\sigma_r^2\mathbf I_{M+1})$ and $\mathbf W_k=[\mathbf w_{k,1},\dots,\mathbf w_{k,M},\mathbf w_{k,r}]$. The echo SNR depends on the beamformers only through $\mathbf R_k$:
\begin{equation}
\rho_{k,n}=\frac{N_s|\alpha_{k,n}|^2\,\mathbf a_{k,n}^H\mathbf R_k\mathbf a_{k,n}}{\sigma_r^2},
\end{equation}
where $\mathbf a_{k,n}\triangleq\mathbf a(\theta_{k,n},\phi_{k,n})$. This is the entry point of the ISAC coupling: sensing performance is invariant to how power is split among streams and depends only on the aggregate spatial covariance, so communication streams also contribute to sensing.

\textbf{Fisher information matrix (FIM).} For each target, the unknown parameter vector is $\boldsymbol\xi_n=[x_n,\,y_n,\,|\alpha_n|,\,\arg\alpha_n]^T$, where the last two entries are nuisance parameters. The FIM consists of two parts.

\emph{Angle/amplitude-phase part.} Define
\begin{multline}
\mathbf H_{k,n}=\big[\alpha_{k,n}\partial_x\mathbf a_{k,n},\ \alpha_{k,n}\partial_y\mathbf a_{k,n},\\
|\alpha_{k,n}|\mathbf a_{k,n},\ j|\alpha_{k,n}|\mathbf a_{k,n}\big]\in\mathbb C^{N_t\times4},
\end{multline}
then
\begin{equation}
\mathbf J^{\rm ang}_n=\frac{2N_s}{\sigma_r^2}\sum_{k}\operatorname{Re}\big\{\mathbf H_{k,n}^H\mathbf R_k\mathbf H_{k,n}\big\}.
\end{equation}

\emph{Time-delay part.} The round-trip delay is $\tau_{k,n}=2d_{k,n}/c$. Assuming a flat spectrum over bandwidth $B$ (effective bandwidth $B_{\rm rms}^2=B^2/12$), the delay-induced information on the horizontal position is
\begin{equation}
\mathbf J^{\rm del}_n=\sum_{k}8\pi^2B_{\rm rms}^2\Big(\frac{2}{c}\Big)^2\rho_{k,n}\,\mathbf g_{k,n}\mathbf g_{k,n}^T,
\end{equation}
with $\mathbf g_{k,n}=(\mathbf q_n-\mathbf p_k)_{xy}/d_{k,n}$, so that
\begin{equation}
\mathbf J_n=\mathbf J^{\rm ang}_n+\operatorname{blkdiag}\big(\mathbf J^{\rm del}_n,\ \mathbf 0_{2\times2}\big).
\end{equation}

\textbf{Sensing metric.} Eliminating the nuisance parameters via the Schur complement, the position CRB of target $n$ is
\begin{equation}
\mathrm{CRB}_n\big(\{\mathbf p_k,\mathbf R_k\}\big)=\big[\mathbf J_n^{-1}\big]_{1,1}+\big[\mathbf J_n^{-1}\big]_{2,2}\ \ [\mathrm{m}^2].
\end{equation}

\textbf{Remark (Fusion assumption).} Summing the angle term over $k$ before inversion implies that all UAVs share a single pair of nuisance parameters, i.e., coherent fusion. This is an optimistic approximation, since the true $\alpha_{k,n}$ differ across UAVs.

\textbf{Remark (Cone ambiguity).} A ULA resolves only the cone angle. When $\Delta y=0$, the angular gradient becomes parallel to the range gradient, $\mathbf J_n$ becomes singular, and the sign of $y$ cannot be resolved ($\mathrm{CRB}\to\infty$). To avoid numerical divergence, we use the capped metric $\mathrm{CRB}^{\rm c}_n=\min\{\mathrm{CRB}_n,C_{\rm cap}\}$. The ambiguity can be removed by a uniform planar array (UPA) or by introducing array yaw as an additional degree of freedom, which is left for future work.

\subsection{UAV Mobility and Safety Constraints}
The trajectory of each UAV is subject to kinematic and safety constraints:
\begin{align}
&\|\mathbf p_k[t+1]-\mathbf p_k[t]\|\le v_{\max}\Delta t,\\
&z_{\min}\le z_k[t]\le z_{\max},\quad (x_k[t],y_k[t])\in\mathcal A,\\
&\|\mathbf p_i[t]-\mathbf p_j[t]\|\ge d_{\min},\quad\forall\,i\neq j,
\end{align}
where $\mathcal A$ is the horizontal mission area and $d_{\min}$ is the inter-UAV collision-avoidance distance.

\begin{figure}[!t]

\centering

\includegraphics[width=\linewidth]{UNNC.png}

\caption{Illustration of the considered system.}

\label{fig:system}

\end{figure}

\subsection{Assumptions}
\textbf{Assumptions.} The model relies on the following assumptions, each with its justification.
\begin{itemize}
  \item \textbf{(A1) Ideal backhaul and global CSI.} Backhaul links are error-free and delay-free, since this isolates the ISAC design from control-plane limitations and yields a performance upper bound.
  \item \textbf{(A2) Perfect CSI and position knowledge.} UAV and user positions, and the channels derived from them, are known exactly, because UAV positioning (GNSS/RTK) and pilot-based estimation are accurate relative to the geometry scales considered.
  \item \textbf{(A3) Block fading.} Channels are constant within a slot and independent across slots, so that one slot corresponds to one beamforming decision and one rate evaluation.
  \item \textbf{(A4) Expected-channel LoS/NLoS model.} The LoS probability enters through an averaged large-scale gain, because it keeps the channel a smooth function of the UAV position, which is required for gradient-based learning.
  \item \textbf{(A5) Spatial degrees of freedom.} $M+1\le N_t$, so that each UAV can transmit one stream per user plus one sensing stream without rank limitation.
  \item \textbf{(A6) Symbol statistics.} Data and probing symbols are zero-mean, unit-power, circularly symmetric, and mutually uncorrelated, because user data are independent and the probing waveform is a deterministic reference, which makes the transmit covariance a sum of beam outer products.
  \item \textbf{(A7) Separated sensing and communication beams.} The sensing stream is dedicated and the communication streams are not waveform-shared, since this decouples the sensing metric from the data content and exposes a clean power-split tradeoff.
  \item \textbf{(A8) Static targets with known altitude.} Targets do not move during the mission and lie on the ground plane, so only the horizontal position needs to be estimated.
  \item \textbf{(A9) Single-channel echo reception with perfect transmit--receive isolation.} Each UAV receives with one channel and suffers no self-interference, because this gives a conservative sensing SNR (no receive-array gain) while avoiding the modeling of full-duplex residual interference.
  \item \textbf{(A10) Coherent fusion across UAVs.} Echo observations of different UAVs are fused with shared nuisance parameters, which gives an optimistic but tractable CRB for cooperative localization.
  \item \textbf{(A11) Kinematics.} UAVs move with bounded speed $v_{\max}$, fixed altitude band, and a minimum separation $d_{\min}$; propulsion energy and array yaw are not modeled, because the focus is on the ISAC resource tradeoff rather than on flight energy.
\end{itemize}

\section{Problem Formulation}
State the task to be solved (e.g., the optimization problem), e.g.,
\begin{subequations}\label{opt:main}
\begin{align}
\min_{\mathbf{x}} \quad & f_0(\mathbf{x}) \label{opt:obj}\\
\text{s.t.}\quad
& f_i(\mathbf{x}) \le 0, \quad i = 1, \ldots, m, \label{opt:ineq}\\
& h_j(\mathbf{x}) = 0, \quad j = 1, \ldots, p. \label{opt:eq}
\end{align}
\end{subequations}
Explain each line of problem~\eqref{opt:main} in words: the meaning of the objective~\eqref{opt:obj}, the physical origin of the inequality constraints~\eqref{opt:ineq}, and why the equality constraints~\eqref{opt:eq} couple the decision variables.

Discuss the challenges of problem~\eqref{opt:main} that motivate the proposed method: non-convexity of $f_0(\cdot)$, coupling across blocks, lack of global state information, or the cost of each iteration. These challenges should map one-to-one onto the design choices of the next section.

\section{Proposed Method}\label{sec:method}
Describe the approach to solving problem~\eqref{opt:main}. The overall procedure is summarized in Algorithm~\ref{alg:proposed}.

\subsection{Overview}
Give the intuition first: which challenge from the previous section each building block addresses, and how the blocks fit together. A reader who stops after this subsection should still understand the main idea of the paper.

\subsection{Iterative Update}
Detail one iteration of Algorithm~\ref{alg:proposed}. With the other blocks fixed, the update of block $k$ reduces to a subproblem of the form
\begin{equation}\label{eq:update}
\mathbf{x}_{k}^{(t+1)} \in \arg\min_{\mathbf{x}_{k}}\; f_0\big(\mathbf{x}_{k}, \mathbf{x}_{-k}^{(t)}\big),
\end{equation}
which admits a closed-form or low-complexity solution because ... Repeat the argument for the remaining blocks and state the cost of one full iteration.

\begin{algorithm}[!t]
\caption{Proposed Algorithm}
\label{alg:proposed}
\begin{algorithmic}[1]
\STATE Initialize $\mathbf{x}^{(0)}$ and set $t = 0$
\REPEAT
\STATE Update one block of $\mathbf{x}$ with the other blocks fixed
\STATE $t \leftarrow t + 1$
\UNTIL convergence in objective~\eqref{opt:obj}
\STATE \textbf{return} $\mathbf{x}^{(t)}$
\end{algorithmic}
\end{algorithm}

\subsection{Convergence and Complexity}
State the convergence guarantee of Algorithm~\ref{alg:proposed}, e.g., a monotone decrease of the objective~\eqref{opt:obj} and convergence to a stationary point under suitable conditions, and bound the overall complexity in terms of the problem dimensions and the number of iterations. Formal proofs are deferred to an appendix or omitted for brevity.

\section{Simulation Results}\label{sec:results}

\subsection{Simulation Setup}
Describe the experimental setup: parameters, baselines, and evaluation metrics. Specify the simulator or dataset, the parameter ranges swept, the baselines chosen together with their citations, and the hardware and runtime environment, so that the results can be reproduced. Unless stated otherwise, each point is averaged over enough random seeds to keep the confidence intervals within the line width of the plots.

\subsection{Numerical Results}
Present the results and main findings. Table~\ref{tab:comparison} gives an example of a three-point IEEE-style table.

Lead the reader through the comparison: which observation matters most, why the proposed scheme wins (linking back to the design choices in Section~\ref{sec:method}), and in which regime the gain shrinks or the assumptions break down. Finish with the take-away message of the evaluation.

\begin{table}[!t]
\caption{An Example Table}
\label{tab:comparison}
\centering
\begin{tabular}{@{}lccc@{}}
\toprule
Scheme & Metric 1 & Metric 2 & Metric 3\\
\midrule
Baseline A & -- & -- & --\\
Baseline B & -- & -- & --\\
Proposed    & -- & -- & --\\
\bottomrule
\end{tabular}
\end{table}

\section{Conclusion}\label{sec:conclusion}
Conclude the paper by restating the problem, the proposed approach, and the headline numerical gain in two or three sentences. Then outline possible future work, e.g., extensions to ... , theoretical analysis of ... , and practical deployment aspects.

\section*{Acknowledgment}
The authors would like to thank...

\bibliographystyle{IEEEtran}
\bibliography{references}

\end{document}
